\documentclass[10pt, aspectratio=169]{beamer}

% --- TEMA MODERNO ---
\usetheme{metropolis} 

% --- COLORI PERSONALIZZATI (Abbinati ai tuoi schemi TikZ) ---
\definecolor{myteal}{HTML}{008080}
\setbeamercolor{palette primary}{bg=myteal, fg=white}
\setbeamercolor{progress bar}{fg=orange, bg=myteal!30}
\setbeamercolor{title separator}{fg=orange}
\setbeamercolor{alerted text}{fg=orange}
\metroset{block=fill} % Aggiunge uno sfondo leggero ai blocchi di testo

\usepackage[utf8]{inputenc}
\usepackage[T1]{fontenc}
\usepackage[italian]{babel}
\usepackage{graphicx}

% --- FONT MODERNO ---
\usepackage[default]{FiraSans} % Font pulito, altamente leggibile e moderno

\usepackage{tikz}
\usetikzlibrary{positioning, arrows.meta, shapes.geometric, calc}

\title{Envass 2.0: un sistema gestionale per l'analisi microbiologica}
\subtitle{Digitalizzazione, Automazione e Integrità del Dato}
\author{Candidato: Christian Mick \\ Relatore: Prof. Guido Sciavicco \\ Correlatore: Dott. Michele Ghiotti}
\institute{Università degli Studi di Ferrara \\ Corso di Laurea in Informatica}
\date{A.A. 2025 - 2026}

\begin{document}
	

	\begin{frame}
		\titlepage
	\end{frame}
	% NUOVA SLIDE 1
	\begin{frame}{Contesto: Evoluzione di ENVASS}
		
		\textbf{CIAS – Centro di Ricerca dell'Università di Ferrara}
		
		\vspace{0.3cm}
		
		ENVASS 2.0 nasce come evoluzione del sistema gestionale esistente
		utilizzato per la gestione e il tracciamento di campioni ambientali
		in ambienti ad alta sterilità.
		
		\vspace{0.5cm}
		
		\begin{columns}[T]
			
			\begin{column}{0.48\textwidth}
				\begin{alertblock}{ENVASS 1.0}
					\begin{itemize}
						\item Funzionalità operative limitate
						\item Processi parzialmente manuali
						\item Minore automazione
						\item Audit e tracciabilità ridotti
					\end{itemize}
				\end{alertblock}
			\end{column}
			
			\begin{column}{0.48\textwidth}
				\begin{block}{ENVASS 2.0}
					\begin{itemize}
						\item Workflow completamente digitalizzato
						\item Generazione automatica dei campioni
						\item Audit centralizzato
						\item Integrità e immutabilità dei dati
					\end{itemize}
				\end{block}
			\end{column}
			
		\end{columns}
		
	\end{frame}
	% NUOVA SLIDE 2
	\begin{frame}{Obiettivi del Progetto}
		
		\begin{columns}[T]
			
			\begin{column}{0.48\textwidth}
				\begin{block}{Esigenze del laboratorio}
					\begin{itemize}
						\item Ridurre errori operativi
						\item Automatizzare attività ripetitive
						\item Garantire tracciabilità completa
						\item Soddisfare requisiti ispettivi
					\end{itemize}
				\end{block}
			\end{column}
			
			\begin{column}{0.48\textwidth}
				\begin{block}{Obiettivi implementativi}
					\begin{itemize}
						\item Modellare il workflow reale
						\item Garantire coerenza del database
						\item Centralizzare audit e logging
						\item Migliorare affidabilità del dato
					\end{itemize}
				\end{block}
			\end{column}
			
		\end{columns}
		
		\vspace{0.6cm}
		
		\centering
		\textbf{
			Dal workflow fisico del laboratorio ad un sistema digitale tracciabile
		}
		
	\end{frame}

\begin{frame}{Stack Tecnologico}
	
	\begin{columns}[T]
		
		% SINISTRA
		\begin{column}{0.42\textwidth}
			
			\begin{block}{Scelte tecnologiche}
				\small
				\begin{itemize}
					\item \textbf{Laravel}\\
					MVC, ORM Eloquent e middleware.
					
					\vspace{0.15cm}
					
					\item \textbf{Bootstrap}\\
					Interfacce responsive e sviluppo rapido.
					
					\vspace{0.15cm}
					
					\item \textbf{jQuery}\\
					Gestione dinamica di eventi e componenti.
					
					\vspace{0.15cm}
					
					\item \textbf{PHP + SQL}\\
					Logica applicativa e persistenza dati.
				\end{itemize}
			\end{block}
			
		\end{column}
		
		% DESTRA
		\begin{column}{0.58\textwidth}
			
			\centering
			
			\resizebox{\textwidth}{!}{
				\begin{tikzpicture}[
					node distance=1cm and 1cm,
					core/.style={draw=orange, thick, fill=orange!10, rectangle, rounded corners, minimum height=1.2cm, minimum width=3.5cm, align=center, font=\Large\bfseries},
					frontend/.style={draw=teal, thick, fill=teal!10, rectangle, rounded corners, minimum height=0.9cm, minimum width=2.8cm, align=center, font=\large\bfseries},
					subfront/.style={draw=teal!60, thick, fill=white, rectangle, rounded corners, minimum height=0.6cm, minimum width=1.5cm, align=center, font=\normalsize},
					backend/.style={draw=gray, thick, fill=gray!20, rectangle, rounded corners, minimum height=1cm, minimum width=3cm, align=center, font=\large\bfseries},
					db/.style={draw=gray, thick, fill=gray!20, cylinder, shape border rotate=90, aspect=0.25, minimum height=1.6cm, minimum width=2.6cm, align=center, font=\large\bfseries},
					line/.style={thick, draw=gray}
					]
					
					\node[core] (laravel) {Laravel 11};
					
					\node[frontend, above=1cm of laravel, xshift=-2.2cm] (bootstrap) {Bootstrap};
					\node[frontend, above=1cm of laravel, xshift=2.2cm] (jquery) {jQuery};
					
					\node[subfront, above=0.5cm of bootstrap, xshift=-1cm] (html) {HTML};
					\node[subfront, above=0.5cm of bootstrap, xshift=1cm] (css) {CSS};
					
					\node[backend, below=1cm of laravel, xshift=-2.2cm] (php) {PHP 8.2};
					\node[db, below=1cm of laravel, xshift=2.2cm] (sql) {SQL};
					
					\draw[line] (bootstrap)--(html);
					\draw[line] (bootstrap)--(css);
					
					\draw[line] (bootstrap)--(laravel);
					\draw[line] (jquery)--(laravel);
					
					\draw[line] (laravel)--(php);
					\draw[line] (php)--(sql);
					
				\end{tikzpicture}
			}
			
		\end{column}
		
	\end{columns}
	
\end{frame}
	

	

	\begin{frame}{La Pipeline Operativa}
		\centering
		\vspace{1cm}
		\resizebox{0.95\textwidth}{!}{
			\begin{tikzpicture}[
				node distance=1.5cm and 2cm,
				box/.style={draw=teal, thick, fill=teal!10, rectangle, rounded corners, minimum height=1.5cm, minimum width=2.5cm, align=center, font=\bfseries}
				]
				% Nodi
				\node[box] (campagna) {1. Campagna \\ di Campionamento};
				\node[box, right=of campagna] (verbali) {2. Verbali \\ (Sopralluogo/Accettazione)};
				\node[box, right=of verbali] (campioni) {3. Campioni \\ (Analisi Microbiologica)};
				\node[box, right=of campioni] (rapporto) {4. Rapporto \\ di Prova (RDP)};
				
				% Frecce
				\draw[-{Latex[length=3mm]}, thick, teal] (campagna) -- (verbali);
				\draw[-{Latex[length=3mm]}, thick, teal] (verbali) -- (campioni);
				\draw[-{Latex[length=3mm]}, thick, teal] (campioni) -- (rapporto);
			\end{tikzpicture}
		}
		\vspace{1.5cm}
		\begin{itemize}
			\item Un flusso sequenziale che mappa esattamente il protocollo fisico del laboratorio.
			\item Ogni fase fornisce i metadati necessari alla fase successiva in modo automatizzato.
		\end{itemize}
	\end{frame}
	

	
	
	\begin{frame}{Generazione Automatica dei Campioni}
		\centering
		\vspace{0.2cm}
		
		\resizebox{0.95\textwidth}{!}{
			\begin{tikzpicture}[
				node distance=0.7cm and 1cm,
				font=\sffamily,
				startblock/.style={draw=teal!70!black, fill=teal!5, thick, rectangle, rounded corners=3pt, minimum height=0.9cm, minimum width=8.5cm, align=center, font=\sffamily\bfseries\color{teal!80!black}},
				datablock/.style={draw=gray!80!black, fill=gray!5, thick, rectangle, rounded corners=3pt, minimum height=1cm, minimum width=6.5cm, align=center},
				actionblock/.style={draw=orange!80!black, fill=orange!5, thick, rectangle, rounded corners=3pt, minimum height=1cm, minimum width=3.8cm, align=center, font=\sffamily\small\bfseries\color{orange!90!black}},
				endblock/.style={draw=teal!70!black, fill=teal!5, thick, rectangle, rounded corners=3pt, minimum height=0.9cm, minimum width=8.5cm, align=center, font=\sffamily\bfseries\color{teal!80!black}},
				arrow/.style={-{Latex[length=2.5mm, width=1.8mm]}, thick, gray!80!black}
				]
				
				\node[startblock] (start) {Convalida del Verbale di Accettazione};
				
				\node[datablock, below=of start] (calc) {Calcolo dei Campioni Attesi \\ \footnotesize (in base a struttura e partizioni)};
				
				\node[actionblock, below left=0.8cm and 0.3cm of calc] (restore) {Ripristino Campioni \\ Eliminati (Soft-Delete)};
				\node[actionblock, below right=0.8cm and 0.3cm of calc] (create) {Creazione Campioni \\ Nuovi/Mancanti};
				
				\node[endblock, below=3.2cm of calc] (end) {Allineamento completato: Database coerente con il prelievo reale};
				
				\draw[arrow] (start) -- (calc);
				
				\draw[arrow] (calc.south) -- ++(0,-0.3) -| (restore.north);
				\draw[arrow] (calc.south) -- ++(0,-0.3) -| (create.north);
				
				\draw[arrow] (restore.south) -- ++(0,-0.4) -| ([xshift=-2.2cm]end.north);
				\draw[arrow] (create.south) -- ++(0,-0.4) -| ([xshift=2.2cm]end.north);
			\end{tikzpicture}
		}
		
		
		\begin{columns}[t]
			\begin{column}{0.48\textwidth}
				\textbf{Idempotenza}\\
				\small Il sistema può rieseguire il calcolo senza duplicare i dati nel database.
			\end{column}
			\begin{column}{0.48\textwidth}
				\textbf{Prevenzione errori}\\
				\small Azzera linearmente il rischio di inserimento manuale non corretto.
			\end{column}
		\end{columns}
	\end{frame}

	% 5. Blocco del Rapporto di Prova (Schema)
	\begin{frame}{Immutabilità: Blocco del Rapporto di Prova}
		\centering
		\vspace{0.2cm}
		\resizebox{0.8\textwidth}{!}{
			\begin{tikzpicture}[
				box/.style={draw=gray, thick, fill=white, rectangle, rounded corners, minimum height=1.2cm, minimum width=4cm, align=center},
				locked/.style={draw=red!80, thick, fill=red!10, rectangle, rounded corners, minimum height=1.2cm, minimum width=4cm, align=center, font=\bfseries},
				arrow/.style={-{Latex[length=3mm]}, thick, red!80}
				]
				\node[locked] (rdp) {Rapporto di Prova (RDP) \\ \textcolor{red!80}{\textbf{[APPROVATO E BLOCCATO]}}};
				
				\node[locked, below left=2cm and 1cm of rdp] (campioni) {Campioni Microbiologici \\ \textcolor{red!80}{\textbf{[IN SOLA LETTURA]}}};
				\node[locked, below right=2cm and 1cm of rdp] (verbali) {Verbali Collegati \\ \textcolor{red!80}{\textbf{[IN SOLA LETTURA]}}};
				
				\draw[arrow] (rdp) -- (campioni) node[midway, left, xshift=-0.2cm, text=black, font=\small, align=center] {Propagazione \\ Blocco};
				\draw[arrow] (rdp) -- (verbali) node[midway, right, xshift=0.2cm, text=black, font=\small, align=center] {Propagazione \\ Blocco};
			\end{tikzpicture}
		}
		\vspace{0.5cm}
		\begin{itemize}
			\item \textbf{Integrità garantita}: L'approvazione finale congela l'intero ramo dei dati pregressi.
			\item In caso di errore, si ricorre alla \textbf{Copia Correttiva}: storicizza il documento errato e genera una nuova revisione tracciata.
		\end{itemize}
	\end{frame}
	\begin{frame}{Immutabilità: Blocco del Rapporto di Prova}
		\centering
		\vspace{0.2cm}
		
		\resizebox{0.85\textwidth}{!}{
			\begin{tikzpicture}[
				node distance=0.8cm and 1.2cm,
				font=\sffamily,
				% Palette cromatica basata su tonalità di rosso desaturato e grigio tecnico
				rootblock/.style={draw=red!80!black, fill=red!5, thick, rectangle, rounded corners=3pt, minimum height=1.2cm, minimum width=5.5cm, align=center},
				leafblock/.style={draw=red!60!black, fill=gray!5, thick, rectangle, rounded corners=3pt, minimum height=1.2cm, minimum width=4.8cm, align=center},
				arrow/.style={-{Latex[length=2.5mm, width=1.8mm]}, thick, red!70!black}
				]
				
				% Nodo principale (Sorgente del blocco)
				\node[rootblock] (rdp) {\textbf{Rapporto di Prova (RDP)} \\ \footnotesize\color{red!80!black}{\textbf{[APPROVATO E BLOCCATO]}}};
				
				% Nodi subordinati (Destinatari del blocco)
				\node[leafblock, below left=1cm and 0.2cm of rdp] (campioni) {\textbf{Campioni Microbiologici} \\ \footnotesize\color{red!70!black}{\textbf{[IN SOLA LETTURA]}}};
				\node[leafblock, below right=1cm and 0.2cm of rdp] (verbali) {\textbf{Verbali Collegati} \\ \footnotesize\color{red!70!black}{\textbf{[IN SOLA LETTURA]}}};
				
				% Punto di derivazione ortogonale per la propagazione
				\coordinate (split) at ($(rdp.south)+(0,-0.4)$);
				
				% Linea di giunzione principale con etichetta unica centrata
				\draw[thick, red!70!black] (rdp.south) -- (split) node[midway, right, font=\footnotesize\sffamily\color{black!80}] {Propagazione Blocco};
				
				% Ramificazioni con frecce direzionali
				\draw[arrow] (split) -| (campioni.north);
				\draw[arrow] (split) -| (verbali.north);
			\end{tikzpicture}
		}
		
		\vspace{0.6cm}
		
		% Organizzazione dei contenuti testuali in colonne bilanciate
		\begin{columns}[t]
			\begin{column}{0.48\textwidth}
				\textbf{Integrità garantita}\\
				\small L'approvazione finale congela l'intero ramo dei dati pregressi, impedendo alterazioni retroattive.
			\end{column}
			\begin{column}{0.48\textwidth}
				\textbf{Copia Correttiva}\\
				\small In caso di errore, il sistema storicizza il documento errato e genera una nuova revisione tracciata.
			\end{column}
		\end{columns}
	\end{frame}
	% 6. Logging 1: Architettura Listener e Observer (Schema)
	\begin{frame}{Audit Logging: Architettura del Sistema}
		\centering
		\resizebox{0.9\textwidth}{!}{
			\begin{tikzpicture}[
				node distance=1cm and 2cm,
				box/.style={draw=teal, thick, fill=white, rectangle, rounded corners, minimum height=1.2cm, minimum width=3cm, align=center},
				sys/.style={draw=orange, thick, fill=orange!10, rectangle, minimum height=1.2cm, minimum width=3cm, align=center, font=\bfseries},
				db/.style={draw=gray, thick, fill=gray!20, cylinder, shape border rotate=90, aspect=0.25, minimum height=2cm, minimum width=2.5cm, align=center},
				arrow/.style={-{Latex[length=3mm]}, thick, gray}
				]
				
				\node[box] (auth) {Eventi di Accesso \\ (Login, Logout, Failed)};
				\node[box, below=1.5cm of auth] (crud) {Operazioni Dati \\ (Create, Update, Delete)};
				
				\node[sys, right=of auth] (listener) {Login Audit \\ Subscriber};
				\node[sys, right=of crud] (observer) {Audit \\ Observer};
				
				\node[sys, fill=teal!10, draw=teal, right=2cm of observer, yshift=1.35cm] (service) {Audit Service \\ (Centralizzazione Payload)};
				
				\node[db, right=of service] (database) {Database \\ Logs};
				
				\draw[arrow] (auth) -- (listener);
				\draw[arrow] (crud) -- (observer);
				\draw[arrow] (listener) -| (service);
				\draw[arrow] (observer) -| (service);
				\draw[arrow] (service) -- (database);
			\end{tikzpicture}
		}
		\vspace{0.5cm}
		\begin{itemize}
			\item Logica \textbf{completamente disaccoppiata} dal codice operativo (Inversione del controllo).
			\item Intercetta automaticamente \textbf{CHI} ha agito, \textbf{COSA} ha cambiato e \textbf{QUANDO}.
		\end{itemize}
	\end{frame}
	
	% 7. Logging 2: Liste + Login
	\begin{frame}{Audit Logging: Liste e Accessi}
		\begin{columns}[T] % L'opzione [T] allinea le colonne in alto
			\begin{column}{0.48\textwidth}
				\begin{block}{Sicurezza e Accessi (Login)}
					Monitoraggio capillare per prevenire e rintracciare anomalie:
					\begin{itemize}
						\item Tracciamento accessi riusciti, logout e \textbf{tentativi falliti}.
						\item Memorizzazione \textbf{Indirizzo IP} e \textbf{User Agent}.
						\item Eventi legati ad account sospesi o disattivati.
					\end{itemize}
				\end{block}
			\end{column}
			\begin{column}{0.48\textwidth}
				\begin{block}{Anagrafiche e Liste}
					Dizionari di base essenziali per la calibrazione del sistema:
					\begin{itemize}
						\item Storico delle modifiche a Metodi e Strumenti.
						\item Prevenzione modifiche silenziose ai parametri normativi.
						\item Registrazione delta (solo i campi alterati).
					\end{itemize}
				\end{block}
			\end{column}
		\end{columns}
	\end{frame}
	
	% 8. Logging 3: Verbali, Campioni, RDP
	\begin{frame}{Audit Logging: Workflow Analitico}
		\textbf{Tracciamento di Verbali, Campioni e Rapporti di Prova}
		\vspace{0.4cm}
		\begin{itemize}
			\item \textbf{Isolamento Delta}: Memorizzazione mirata. Il sistema calcola l'intersezione tra i dati pre e post-modifica. Vengono salvati \textit{solo} i campi variati.
			\item \textbf{Tracciamento Gerarchico (\texttt{root\_id})}: Le modifiche ai singoli campioni mantengono un legame logico con il documento radice (Verbale o RDP), permettendo estrazioni rapide.
			\item \textbf{Filtro Metadati}: Esclusione automatica nei log di chiavi tecniche o password per favorire la leggibilità ispettiva.
		\end{itemize}
		\vspace{0.3cm}
		\begin{alertblock}{Esito per gli Ispettori}
			Una cronologia trasparente, immutabile e immediatamente consultabile che rispetta al 100\% i criteri di qualità ACCREDIA.
		\end{alertblock}
	\end{frame}
	
	% 9. Grazie
	\begin{frame}[standout]
		Grazie per l'attenzione!
	\end{frame}
	
	% 10. Bibliografia
	\begin{frame}{Bibliografia}
		\begin{thebibliography}{9}
			\bibitem{html} HTML Standard - WHATWG.
			\bibitem{robbins} J. N. Robbins, \textit{Learning Web Design: A Beginner's Guide}, 6th ed., O'Reilly, 2025.
			\bibitem{ecma} ECMA-262, 16th Edition, ECMAScript Language Specification, 2025.
			\bibitem{laravel} Laravel Documentation (v11/13.x) - \texttt{laravel.com/docs}
			\bibitem{atzeni} P. Atzeni, \textit{Basi di dati: modelli e linguaggi di interrogazione}, McGraw-Hill, 2009.
			\bibitem{who} WHO Bacterial Priority Pathogens List 2024, World Health Organization.
		\end{thebibliography}
	\end{frame}
	
	\begin{frame}
		\titlepage
	\end{frame}
	
\end{document}